\usepackage{amsthm}
\newtheorem{theorem}{Theorem}[section]
\newtheorem{definition}[theorem]{Definition}
\newtheorem{corollary}[theorem]{Corollary}
\newtheorem{lemma}[theorem]{Lemma}
\begin{document}
\section{Introduction}
Here is some text to introduce the theorems, definitions, corollaries, and lemmas.
\section{Main Result}
\begin{theorem}
This is the Pythagorean theorem.
\label{pythagorean}
This theorem states that in a right triangle:
\[x^2 + y^2 = z^2\]
\end{theorem}
\begin{proof}
This is the proof of the theorem.
\end{proof}
\begin{definition}[Fibration]
A fibration is a mapping between two topological spaces that has the homotopy lifting
property for every space \(X\).
\end{definition}
\begin{corollary}
There is no right triangle whose sides measure 3 cm, 4 cm, and 6 cm.
\end{corollary}
\begin{lemma}
Given two line segments whose lengths are \(a\) and \(b\), there exists a real number \(r\) such
that \(b = ra\).
\end{lemma}
\section{Discussion}
Some discussion about the results.
\end{document}
